\section*{Modèle retenu : application au CO$_2$ et CO}

Dans un premier temps, on retient le modèle réduit obtenu à partir de
l'approche de Thiele. Le transport est supposé purement diffusif, le
système est considéré en régime stationnaire, et les effets de migration,
de convection et de chute ohmique sont négligés.

On considère uniquement deux espèces : le CO$_2$, consommé par la
réaction, et le CO, produit par cette même réaction. La réaction globale
est représentée sous la forme simplifiée :

\begin{equation}
\mathrm{CO_2 \longrightarrow CO}.
\end{equation}

Cette écriture ne décrit pas explicitement les espèces ioniques, l'eau,
ni les intermédiaires réactionnels. Elle sert uniquement à construire un
premier modèle diffusion--réaction pour les deux espèces carbonées.

Le système étudié est donc :

\begin{equation}
\begin{cases}
\varepsilon_l D_{CO_2}^{\mathrm{eff}}
\dfrac{d^2 C_{CO_2}}{dz^2}
=
a_v k C_{CO_2},\\[0.4cm]
\varepsilon_l D_{CO}^{\mathrm{eff}}
\dfrac{d^2 C_{CO}}{dz^2}
=
- a_v k C_{CO_2}.
\end{cases}
\end{equation}

On introduit la coordonnée adimensionnée

\begin{equation}
\zeta = \frac{z}{l_e},
\end{equation}

ainsi que le module de Thiele associé à la consommation du CO$_2$ :

\begin{equation}
\phi^2 =
\frac{a_v k l_e^2}
{\varepsilon_l D_{CO_2}^{\mathrm{eff}}}.
\end{equation}

Le problème devient alors :

\begin{equation}
\begin{cases}
\dfrac{d^2 C_{CO_2}}{d\zeta^2}
=
\phi^2 C_{CO_2},\\[0.4cm]
\dfrac{d^2 C_{CO}}{d\zeta^2}
=
-\dfrac{a_v k l_e^2}
{\varepsilon_l D_{CO}^{\mathrm{eff}}}
C_{CO_2}.
\end{cases}
\end{equation}

On impose une concentration fixée à l'interface avec le bulk et un flux
nul à la paroi :

\begin{align}
C_{CO_2}(0) &= C_{CO_2}^{\mathrm{bulk}},
&
\frac{d C_{CO_2}}{d\zeta}(-1) &= 0,\\
C_{CO}(0) &= C_{CO}^{\mathrm{bulk}},
&
\frac{d C_{CO}}{d\zeta}(-1) &= 0.
\end{align}

La solution analytique pour le CO$_2$ est :

\begin{equation}
\boxed{
C_{CO_2}(\zeta)
=
C_{CO_2}^{\mathrm{bulk}}
\frac{
\cosh\!\left[\phi(\zeta+1)\right]
}
{\cosh(\phi)}
}
\end{equation}

Le CO étant produit localement à partir de la consommation du CO$_2$,
son profil est obtenu en injectant la solution précédente dans l'équation
de conservation du CO. On obtient :

\begin{equation}
\boxed{
C_{CO}(\zeta)
=
C_{CO}^{\mathrm{bulk}}
+
\frac{D_{CO_2}^{\mathrm{eff}}}
     {D_{CO}^{\mathrm{eff}}}
C_{CO_2}^{\mathrm{bulk}}
\left[
1-
\frac{
\cosh\!\left[\phi(\zeta+1)\right]
}
{\cosh(\phi)}
\right]
}
\end{equation}

Ce modèle constitue le premier cas test retenu pour les simulations.
Il permet d'étudier l'influence du module de Thiele sur la pénétration
du CO$_2$ dans l'électrode et sur la formation locale de CO.